% ============================================================
%  RAPPORT DE STAGE / PROJET – AfriMarket
%  Fichier principal : rapport_AfriMarket.tex
%  Compilé avec : pdflatex (2 passes) puis bibtex puis 2 passes
%  Encodage : UTF-8
% ============================================================

\documentclass[12pt, a4paper, oneside]{report}

% ─── Encodage et langue ────────────────────────────────────
\usepackage[utf8]{inputenc}
\usepackage[T1]{fontenc}
\usepackage[french]{babel}

% ─── Mise en page ──────────────────────────────────────────
\usepackage[
  top=2.5cm, bottom=2.5cm,
  left=3cm,  right=2.5cm
]{geometry}
\usepackage{setspace}
\onehalfspacing                    % interligne 1.5

% ─── Typographie ───────────────────────────────────────────
\usepackage{lmodern}               % polices vectorielles
\usepackage{microtype}             % amélioration micro-typographique
\usepackage{csquotes}              % guillemets corrects en français

% ─── Couleurs et graphiques ────────────────────────────────
\usepackage[dvipsnames,svgnames,x11names]{xcolor}
\usepackage{graphicx}
\usepackage{tikz}
\usetikzlibrary{shapes.geometric, arrows.meta, positioning, fit, backgrounds}
\usepackage{pgfplots}
\pgfplotsset{compat=1.18}

% ─── Tableaux ──────────────────────────────────────────────
\usepackage{booktabs}              % filets professionnels
\usepackage{tabularx}              % tableaux à largeur fixe
\usepackage{multirow}              % fusion de lignes
\usepackage{longtable}             % tableaux multi-pages
\usepackage{array}                 % colonnes personnalisées
% Bordures complètes sur tous les tableaux via \hline
\renewcommand{\arraystretch}{1.3}  % hauteur de ligne agréable

% ─── Listes et numérotation ────────────────────────────────
\usepackage{enumitem}
\usepackage{pifont}

% ─── Figures et légendes ───────────────────────────────────
\usepackage{caption}
\usepackage{subcaption}
\usepackage{float}                 % option [H] pour les figures
\captionsetup{
  labelfont=bf,
  textfont=it,
  justification=centering
}

% ─── Code source ───────────────────────────────────────────
\usepackage{listings}
\lstset{
  basicstyle=\ttfamily\small,
  keywordstyle=\color{NavyBlue}\bfseries,
  commentstyle=\color{Gray}\itshape,
  stringstyle=\color{ForestGreen},
  numbers=left, numberstyle=\tiny\color{Gray},
  breaklines=true,
  frame=single, rulecolor=\color{Gray!40},
  backgroundcolor=\color{Gray!5},
  tabsize=2
}

% ─── Liens hypertextes ─────────────────────────────────────
\usepackage[
  colorlinks=true,
  linkcolor=NavyBlue,
  citecolor=BrickRed,
  urlcolor=Teal,
  bookmarks=true,
  pdfauthor={[Votre Prénom NOM]},
  pdftitle={Rapport de projet – AfriMarket}
]{hyperref}

% ─── Bibliographie ─────────────────────────────────────────
\usepackage[
  backend=bibtex,
  style=numeric,
  sorting=nyt
]{biblatex}
\addbibresource{references.bib}

% ─── Abréviations (acronymes) ──────────────────────────────
\usepackage[
  acronym,
  toc,
  nonumberlist,
  nopostdot
]{glossaries}
\makeglossaries

% ─── En-têtes et pieds de page ─────────────────────────────
\usepackage{fancyhdr}
\pagestyle{fancy}
\fancyhf{}
\fancyhead[L]{\small\leftmark}
\fancyhead[R]{\small AfriMarket}
\fancyfoot[C]{\thepage}
\renewcommand{\headrulewidth}{0.4pt}

% ─── Numérotation des sections dans le sommaire ────────────
\setcounter{tocdepth}{3}
\setcounter{secnumdepth}{3}

% ─── Formatage des chapitres ───────────────────────────────
\usepackage{titlesec}
\titleformat{\chapter}[block]
  {\normalfont\Large\bfseries\color{NavyBlue}}
  {\thechapter.}{1em}{}
  [\vspace{0.2em}\titlerule]

% ─── Commandes utiles ──────────────────────────────────────
\newcommand{\TODO}[1]{\textcolor{red}{\textbf{[TODO: #1]}}}
% Commande pour mettre en avant un terme technique
\newcommand{\tech}[1]{\texttt{#1}}

% ============================================================
%  DÉFINITION DES ACRONYMES
% ============================================================
\newacronym{api}{API}{Application Programming Interface}
\newacronym{rest}{REST}{Representational State Transfer}
\newacronym{jwt}{JWT}{JSON Web Token}
\newacronym{ui}{UI}{User Interface}
\newacronym{ux}{UX}{User Experience}
\newacronym{mvvm}{MVVM}{Model-View-ViewModel}
\newacronym{p2p}{P2P}{Peer-to-Peer}
\newacronym{crud}{CRUD}{Create Read Update Delete}
\newacronym{jpa}{JPA}{Jakarta Persistence API}
\newacronym{sql}{SQL}{Structured Query Language}
\newacronym{momo}{MoMo}{Mobile Money}
\newacronym{json}{JSON}{JavaScript Object Notation}
\newacronym{spa}{SPA}{Single Page Application}
\newacronym{tsx}{TSX}{TypeScript XML}
\newacronym{orm}{ORM}{Object-Relational Mapping}
\newacronym{http}{HTTP}{HyperText Transfer Protocol}
\newacronym{fcfa}{FCFA}{Franc CFA}
\newacronym{uml}{UML}{Unified Modeling Language}

% ============================================================
%  DÉBUT DU DOCUMENT
% ============================================================
\begin{document}

% ─── Page de titre ─────────────────────────────────────────
\begin{titlepage}
  \centering
  % Logos côte à côte : UIGEE à gauche, Keyce à droite
  % Remplacez les chemins par vos fichiers réels
  \begin{minipage}[c]{0.35\textwidth}
    \centering
    \fbox{\rule{0pt}{1.5cm}\rule{3cm}{0pt}}\\[2pt]
    % \includegraphics[width=3cm]{logos/uigee.png}  % ← décommenter + placer logo
    {\small UIGEE}
  \end{minipage}
  \hfill
  \begin{minipage}[c]{0.35\textwidth}
    \centering
    \fbox{\rule{0pt}{1.5cm}\rule{3cm}{0pt}}\\[2pt]
    % \includegraphics[width=3cm]{logos/keyce.png}  % ← décommenter + placer logo
    {\small Keyce Informatique \& IA}
  \end{minipage}

  \vspace{1.5cm}
  {\small Université / Institut de l'Ingénierie et des Grandes Écoles}\\[0.3cm]
  {\small En partenariat avec Keyce Informatique \& IA}

  \vspace{1cm}
  \rule{\linewidth}{1pt}\\[0.4cm]
  {\LARGE\bfseries\color{NavyBlue}
    Rapport de Projet de Développement\\[0.3cm]
    Application Mobile \textit{AfriMarket}\\[0.2cm]
    Marketplace Agricole P2P pour l'Afrique Centrale
  }\\[0.3cm]
  \rule{\linewidth}{1pt}

  \vspace{1cm}
  {\large Présenté par :}\\[0.4cm]
  {\Large\bfseries [Votre Prénom NOM]}\\[0.2cm]
  {\large Matricule : \textbf{[XXXX-XXXXXX]}}

  \vspace{1cm}
  \begin{tabular}{ll}
    \textbf{Filière :}       & Développement Mobile / Génie Logiciel \\[4pt]
    \textbf{Niveau :}        & [Licence / Master] \\[4pt]
    \textbf{Encadreur :}     & [Nom de l'encadreur] \\[4pt]
    \textbf{Année académique :} & 2024 -- 2025 \\
  \end{tabular}

  \vfill
  {\small Soutenu le \textbf{[JJ Mois AAAA]} devant le jury composé de :}\\[0.3cm]
  \begin{tabular}{ll}
    Président  & [Nom] \\
    Rapporteur & [Nom] \\
    Examinateur & [Nom] \\
  \end{tabular}
\end{titlepage}

% ─── Pages liminaires : numérotation romaine ───────────────
\pagenumbering{roman}
\setcounter{page}{1}

% ============================================================
%  REMERCIEMENTS
% ============================================================
\chapter*{Remerciements}
\addcontentsline{toc}{chapter}{Remerciements}

Nous adressons en premier lieu nos plus sincères remerciements aux membres du jury qui ont
bien voulu consacrer de leur précieux temps à l'évaluation de ce travail. Leurs remarques
constructives contribueront à l'amélioration de nos pratiques professionnelles.

Nos remerciements vont ensuite à \textbf{[Nom du professeur encadreur]}, enseignant
référent, dont la disponibilité, les orientations pédagogiques et les conseils avisés ont
guidé chaque étape de ce projet. Sa rigueur académique nous a permis d'aborder les
difficultés techniques avec méthode.

Nous exprimons également notre gratitude à l'ensemble des tuteurs et maîtres de stage qui
ont suivi l'évolution de ce travail au quotidien, pour leur patience et la transmission
de leur expérience professionnelle.

Enfin, nous remercions les membres de l'administration de l'UIGEE et de Keyce Informatique
\& IA pour leur soutien institutionnel et les conditions d'étude qu'ils ont su créer.

% ============================================================
%  RÉSUMÉ (français)
% ============================================================
\chapter*{Résumé}
\addcontentsline{toc}{chapter}{Résumé}

AfriMarket est une application mobile de type marketplace agricole pair-à-pair conçue pour
le contexte camerounais et, plus largement, pour l'Afrique centrale francophone. Elle met
en relation directe des producteurs agricoles avec des consommateurs urbains, en s'appuyant
sur un réseau de points relais pour la logistique de livraison et sur le Mobile Money pour
les paiements. Le projet comprend trois composantes logicielles : une application
\emph{consommateur} (AfriMarket Client), une application \emph{producteur} (AfriMarket
Producteur) et une interface d'administration web, toutes alimentées par un serveur backend
RESTful.

Ce rapport décrit le processus complet qui a conduit de l'identification du problème de
distribution agricole en milieu urbain jusqu'à la livraison d'une solution fonctionnelle.
Après une analyse de l'état de l'art des plateformes agri-tech existantes, nous exposons
les choix de conception, l'architecture logicielle retenue et la mise en œuvre technique,
avant de présenter les résultats obtenus et les perspectives d'amélioration.

\vspace{1em}
\noindent\textbf{Mots-clés :} application mobile, marketplace agricole, React Native,
Spring Boot, Mobile Money, Cameroun, pair-à-pair, logistique, \gls{rest}.

% ============================================================
%  ABSTRACT (anglais)
% ============================================================
\chapter*{Abstract}
\addcontentsline{toc}{chapter}{Abstract}

AfriMarket is a peer-to-peer agricultural marketplace mobile application designed for the
Cameroonian context and, more broadly, for French-speaking Central Africa. It directly
connects agricultural producers with urban consumers through a relay-point delivery network
and Mobile Money payments. The project consists of three software components: a
\emph{consumer} application (AfriMarket Client), a \emph{producer} application (AfriMarket
Producteur), and a web-based administration interface, all powered by a RESTful backend
server.

This report describes the complete process from the identification of the agricultural
distribution problem in urban areas to the delivery of a functional solution. After an
analysis of the state of the art in existing agri-tech platforms, we present the design
choices, software architecture, and technical implementation, before discussing results and
future improvements.

\vspace{1em}
\noindent\textbf{Keywords:} mobile application, agricultural marketplace, React Native,
Spring Boot, Mobile Money, Cameroon, peer-to-peer, logistics, REST.

% ============================================================
%  LISTE DES FIGURES (auto)
% ============================================================
\listoffigures
\addcontentsline{toc}{chapter}{Liste des figures}

% ============================================================
%  LISTE DES TABLEAUX (auto)
% ============================================================
\listoftables
\addcontentsline{toc}{chapter}{Liste des tableaux}

% ============================================================
%  ABRÉVIATIONS (auto via glossaries)
% ============================================================
\printglossary[type=\acronymtype, title={Liste des Abréviations}]
\addcontentsline{toc}{chapter}{Liste des Abréviations}

% ============================================================
%  GLOSSAIRE
% ============================================================
\chapter*{Glossaire}
\addcontentsline{toc}{chapter}{Glossaire}

\begin{description}[leftmargin=3.5cm, style=nextline]
  \item[Marketplace]
    Plateforme numérique mettant en relation des offreurs et des demandeurs de biens ou
    de services sans que la plateforme ne soit elle-même propriétaire des biens.
  \item[Mobile Money]
    Service financier permettant d'effectuer des paiements, transferts et retraits d'argent
    via un téléphone mobile, très répandu en Afrique subsaharienne
    (Orange Money, MTN MoMo, etc.).
  \item[Point relais]
    Lieu physique tiers (boutique, dépôt) servant d'intermédiaire de livraison entre le
    producteur et le consommateur, évitant ainsi la livraison à domicile.
  \item[Offre groupée]
    Mécanisme par lequel un producteur ne lance la production ou la livraison que si la
    quantité minimale commandée est atteinte par l'ensemble des acheteurs.
  \item[Wallet virtuel]
    Portefeuille électronique interne à l'application, créédité lors du paiement d'une
    commande et débitable vers le compte Mobile Money du producteur.
  \item[Token JWT]
    Jeton numérique signé permettant d'authentifier un utilisateur auprès de l'API sans
    stocker d'état côté serveur.
  \item[Expo / Expo Router]
    Framework JavaScript facilitant le développement d'applications React Native
    multi-plateformes (iOS, Android, Web) avec un système de routage par fichiers.
  \item[Spring Boot]
    Framework Java permettant de créer des applications web autonomes avec un minimum
    de configuration.
\end{description}

% ============================================================
%  TABLE DES MATIÈRES
% ============================================================
\tableofcontents

% ─── Corps du document : numérotation arabe ────────────────
\pagenumbering{arabic}
\setcounter{page}{1}

% ============================================================
%  CHAPITRE 1 – INTRODUCTION GÉNÉRALE
% ============================================================
\chapter{Introduction générale}
\markboth{Chapitre 1 -- Introduction générale}{}

Le développement de l'économie numérique en Afrique subsaharienne connaît une croissance
soutenue portée par la progression rapide de la téléphonie mobile et des infrastructures
Internet. Dans ce contexte, le secteur agricole, qui représente une part significative du
\gls{pib} de nombreux pays d'Afrique centrale, reste paradoxalement l'un des moins
numérisés. Les agriculteurs manquent d'outils adaptés pour valoriser leur production et
rejoindre efficacement les marchés urbains, tandis que les consommateurs souffrent d'une
opacité persistante sur la provenance et la qualité des denrées alimentaires.

Ce projet s'inscrit dans cette dynamique de transformation numérique de l'agriculture en
proposant AfriMarket, une plateforme mobile pair-à-pair dédiée à la mise en relation
directe de producteurs agricoles et de consommateurs au Cameroun. L'idée centrale est
d'éliminer les intermédiaires traditionnels souvent à l'origine des marges excessives et
des pertes post-récolte, tout en s'appuyant sur des points relais locaux pour gérer la
logistique de livraison et sur le Mobile Money pour sécuriser les paiements.

La problématique centrale de ce travail peut se formuler ainsi : comment concevoir et
développer une application mobile qui s'adapte aux contraintes spécifiques du marché
camerounais — connectivité intermittente, prédominance du paiement mobile, faible culture
numérique de certains producteurs — tout en offrant une expérience utilisateur fluide et
sécurisée à l'ensemble des acteurs de la chaîne ? Cette question en recouvre plusieurs
autres : comment modéliser un système d'offres groupées permettant aux producteurs de ne
livrer que lorsqu'un seuil minimal de commandes est atteint ? Comment orchestrer le cycle
de vie d'une commande entre les rôles de producteur, de consommateur, de gestionnaire de
point relais et d'administrateur ? Comment intégrer un portefeuille virtuel interne tout
en prévoyant les passerelles vers les opérateurs de Mobile Money ?

La motivation principale de ce projet naît d'un constat terrain : au marché central de
Yaoundé comme dans les marchés de quartier, les agriculteurs des zones péri-urbaines
vendent souvent leur production à des prix très inférieurs à ceux affichés en ville, faute
de moyens logistiques et de visibilité. AfriMarket vise à corriger cette asymétrie en
dotant les producteurs d'un outil de gestion d'offres accessible depuis un simple
smartphone et en permettant aux consommateurs de s'organiser collectivement pour atteindre
les seuils d'achat. Au-delà de l'enjeu économique, le projet répond à une ambition sociale :
renforcer la souveraineté alimentaire locale et valoriser le travail agricole.

Les objectifs de ce projet sont multiples et structurés en niveaux de priorité.
L'objectif principal est de livrer une application mobile fonctionnelle comprenant trois
interfaces distinctes — consommateur, producteur et administration — reliées à un backend
\gls{rest} unifié. Les objectifs secondaires incluent la mise en place d'un système
d'offres groupées avec progression visible en temps réel, l'intégration d'un wallet
virtuel avec mécanisme de retrait vers Mobile Money, la gestion d'un réseau de points
relais et la sécurisation des échanges par authentification \gls{jwt}.

Ce rapport est organisé en trois chapitres principaux. Le premier chapitre — cette
introduction — pose le contexte, la problématique, les motivations et les objectifs du
projet. Le deuxième chapitre explore les concepts fondamentaux mobilisés et dresse un
état de l'art des solutions existantes. Le troisième chapitre détaille la méthodologie
suivie : analyse des besoins, conception et architecture. Le quatrième chapitre présente
les choix technologiques et les résultats obtenus à travers des captures d'écran et une
description des fonctionnalités livrées.

% ============================================================
%  CHAPITRE 2 – CONCEPTS ET ÉTAT DE L'ART
% ============================================================
\chapter{Concepts et état de l'art}
\markboth{Chapitre 2 -- Concepts et état de l'art}{}

% ─── 2.1 Concepts fondamentaux ─────────────────────────────
\section{Concepts fondamentaux}

La compréhension d'AfriMarket nécessite la maîtrise de plusieurs concepts transversaux
issus du commerce électronique, de la logistique et du développement mobile.

\subsection{La marketplace pair-à-pair (\glsentryshort{p2p})}

Une marketplace \gls{p2p} est une plateforme numérique qui facilite des transactions
directes entre particuliers ou entre petits producteurs et consommateurs, sans que
l'opérateur de la plateforme ne soit propriétaire des biens échangés. Contrairement aux
places de marché B2C classiques (Business-to-Consumer), le modèle \gls{p2p} repose sur la
confiance entre pairs, renforcée par des mécanismes de notation, de vérification d'identité
et d'escrow (séquestre des fonds). Airbnb et Uber en sont les exemples emblématiques ;
dans le secteur agricole, des plateformes comme Twiga Foods au Kenya ou Hello Tractor en
Afrique de l'Ouest illustrent le potentiel de ce modèle pour le continent africain
\cite{twiga2023}.

\subsection{Les offres groupées (group buying)}

Le mécanisme d'offre groupée, popularisé par Groupon à l'international, consiste à
conditionner l'activation d'une offre commerciale à l'atteinte d'un seuil minimal de
participants. Dans le contexte agricole, ce mécanisme présente un double avantage : il
garantit au producteur une taille critique de commande qui justifie le déplacement ou la
récolte spécifique, et il permet aux consommateurs de bénéficier de prix de gros tout en
achetant en quantité raisonnable. AfriMarket adapte ce concept en affichant un indicateur
de progression en temps réel (pourcentage de la quantité minimale atteinte), ce qui crée
une dynamique sociale d'entraînement entre les acheteurs.

\subsection{Le Mobile Money en Afrique centrale}

Le Mobile Money désigne un service financier permettant d'effectuer des transactions
monétaires via un téléphone mobile, sans nécessiter de compte bancaire traditionnel. En
Afrique subsaharienne, deux opérateurs dominent le marché camerounais : MTN Mobile Money
(MoMo) et Orange Money. Avec un taux de bancarisation inférieur à 20\% au Cameroun
\cite{banquemondiale2023}, le Mobile Money représente le principal vecteur de paiement
numérique pour une large part de la population, y compris les agriculteurs. L'intégration
de ces services dans AfriMarket est donc non seulement souhaitable mais indispensable pour
l'adoption de la plateforme.

\subsection{Le réseau de points relais}

Un point relais est un établissement commercial partenaire (kiosque, boutique, dépôt)
servant d'intermédiaire entre le producteur et le consommateur pour la remise des
commandes. Ce modèle logistique, bien établi en Europe (Mondial Relay, Chronopost), se
révèle particulièrement adapté aux villes africaines où l'adressage postal est souvent
imprécis et où la livraison à domicile génère des coûts importants. AfriMarket intègre un
réseau de points relais géo-référencés, associés à des zones géographiques calquées sur les
arrondissements de Yaoundé.

\subsection{L'architecture RESTful et les microservices}

L'\gls{api} \gls{rest} constitue le style architectural dominant pour les backend
d'applications mobiles. Elle repose sur le protocole \gls{http} et organise les ressources
autour de verbes sémantiques (GET, POST, PUT, DELETE). Dans AfriMarket, cette architecture
permet de découpler complètement les trois interfaces clientes (consommateur, producteur,
admin) du backend, offrant une flexibilité maximale pour l'évolution indépendante de
chaque composante \cite{fielding2000}.

\subsection{La sécurité par JSON Web Token}

Le \gls{jwt} est un standard ouvert (RFC 7519) définissant une manière compacte et
auto-portée de transmettre des informations entre parties sous forme d'objet \gls{json}
signé numériquement. Dans le contexte d'une \gls{api} \gls{rest} sans état (stateless),
le \gls{jwt} remplace avantageusement la session serveur traditionnelle : le token, émis
lors de la connexion et stocké côté client, est joint à chaque requête et vérifié par le
serveur sans consultation d'une base de données de session \cite{jones2015}.

% ─── 2.2 État de l'art ─────────────────────────────────────
\section{État de l'art des solutions existantes}

\subsection{Plateformes agricoles africaines}

Plusieurs initiatives numériques tentent de résoudre le problème de mise en marché des
productions agricoles en Afrique. Twiga Foods (Kenya) propose un modèle B2B qui achète les
productions directement aux agriculteurs pour les revendre aux épiceries urbaines. Hello
Tractor (Nigeria) se concentre sur la location de matériel agricole. Esoko (Ghana) fournit
des informations de prix de marché par SMS. En Afrique centrale, les solutions spécifiques
au Cameroun restent rares et rarement adaptées aux spécificités linguistiques et
logistiques locales.

\begin{table}[H]
  \centering
  \caption{Comparaison des plateformes agri-tech existantes}
  \label{tab:comparaison_plateformes}
  \begin{tabular}{|l|l|l|l|l|}
    \hline
    \textbf{Plateforme} & \textbf{Pays} & \textbf{Modèle} & \textbf{Paiement} & \textbf{Livraison} \\
    \hline
    Twiga Foods   & Kenya     & B2B       & Mobile Money & Propre fleet \\
    \hline
    Hello Tractor & Nigeria   & Location  & Carte/Mobile & Non concerné \\
    \hline
    Esoko         & Ghana     & Info prix & -- & -- \\
    \hline
    Jumia Food    & Multi     & B2C resto & Mobile/Carte & Coursiers \\
    \hline
    \textbf{AfriMarket} & \textbf{Cameroun} & \textbf{P2P} & \textbf{MoMo} & \textbf{Points relais} \\
    \hline
  \end{tabular}
\end{table}

\subsection{Frameworks de développement mobile}

Le choix d'un framework de développement mobile est crucial car il détermine la
maintenance, les performances et la capacité à cibler plusieurs plateformes simultanément.
Trois approches coexistent : le développement natif (Swift/Kotlin), le développement
hybride (Ionic/Cordova) et le développement cross-platform (Flutter, React Native).
React Native, soutenu par Meta, se distingue par son utilisation de JavaScript/TypeScript,
sa large communauté et son écosystème Expo qui simplifie considérablement le processus de
build et de déploiement \cite{eisenman2015}. Il s'impose comme le choix naturel dans un
contexte académique où la réutilisation des compétences web JavaScript est un avantage.

\subsection{Frameworks backend pour les API mobiles}

Spring Boot, maintenu par VMware Pivotal, est le framework Java le plus utilisé en
entreprise pour construire des API RESTful. Sa convention-over-configuration, sa gestion
intégrée de la sécurité via Spring Security, et son écosystème JPA/Hibernate pour la
persistance des données en font un choix robuste et éprouvé. La version 3.x, basée sur
Jakarta EE, apporte des améliorations significatives en termes de performances et de
support des architectures réactives \cite{springboot2024}.

% ─── 2.3 Bilan ─────────────────────────────────────────────
\section{Bilan}

L'analyse de l'existant révèle un écosystème agri-tech africain encore fragmenté, dominé
par des solutions B2B ou des services d'information incomplets. Aucune plateforme ne
propose simultanément les offres groupées, le réseau de points relais géolocalisés, le
wallet virtuel et l'interface tripartite (producteur / consommateur / admin) adaptée au
marché camerounais. AfriMarket comble ce vide en s'appuyant sur des technologies modernes
(React Native avec Expo, Spring Boot 4) et en intégrant nativement les contraintes locales
(Mobile Money, gestion par zones, mode hors-ligne partiel). Ce positionnement justifie
pleinement le développement d'une solution sur mesure plutôt que l'adaptation d'une
solution existante.

% ============================================================
%  CHAPITRE 3 – MÉTHODOLOGIE
% ============================================================
\chapter{Méthodologie}
\markboth{Chapitre 3 -- Méthodologie}{}

% ─── 3.1 Analyse des besoins ───────────────────────────────
\section{Analyse des besoins}

\subsection{Besoins fonctionnels}

Les besoins fonctionnels ont été identifiés en interrogeant les utilisateurs cibles et en
analysant les processus actuels de vente agricole. Ils ont été organisés autour des
cinq acteurs principaux du système.

\subsubsection{Consommateur}

Le consommateur doit pouvoir s'inscrire et se connecter via son numéro de téléphone, parcourir
et filtrer les offres actives par catégorie (légumes, fruits, tubercules, céréales,
élevage) et par zone géographique, visualiser la progression d'une offre groupée, rejoindre
une offre en spécifiant la quantité souhaitée, payer via Mobile Money, sélectionner un
point relais pour la livraison, confirmer la réception de sa commande et ouvrir un litige
en cas de problème.

\subsubsection{Producteur}

Le producteur doit pouvoir créer une offre avec titre, description, catégorie, unité, prix
unitaire, quantité disponible, seuil minimal et date d'expiration, téléverser des photos de
sa production, soumettre l'offre à validation administrative, suivre en temps réel les
commandes reçues, confirmer l'expédition vers les relais, gérer son portefeuille virtuel
et demander un retrait vers son compte Mobile Money.

\subsubsection{Gestionnaire de point relais}

Le gestionnaire de point relais doit recevoir les notifications de dépôt, signaler la
réception des colis et confirmer les remises aux consommateurs.

\subsubsection{Administrateur}

L'administrateur doit valider ou rejeter les offres et les profils de producteurs, accéder
à un tableau de bord avec les indicateurs clés (chiffre d'affaires, commandes, utilisateurs),
gérer les zones, les points relais, les utilisateurs et les paramètres système.

\subsubsection{Finance}

L'opérateur finance doit visualiser l'ensemble des transactions, déclencher les virements
Mobile Money aux producteurs et gérer les remboursements en cas de litige.

\begin{table}[H]
  \centering
  \caption{Matrice des besoins fonctionnels par acteur}
  \label{tab:besoins_fonctionnels}
  \begin{tabularx}{\textwidth}{|l|X|c|}
    \hline
    \textbf{Acteur} & \textbf{Besoin fonctionnel principal} & \textbf{Priorité} \\
    \hline
    Consommateur & Parcourir et rejoindre des offres & Haute \\
    \hline
    Consommateur & Payer via Mobile Money & Haute \\
    \hline
    Consommateur & Confirmer la réception / ouvrir litige & Moyenne \\
    \hline
    Producteur & Créer et publier une offre groupée & Haute \\
    \hline
    Producteur & Gérer son wallet virtuel & Haute \\
    \hline
    Producteur & Retirer ses gains vers MoMo & Haute \\
    \hline
    Relais & Confirmer dépôt et remise & Moyenne \\
    \hline
    Admin & Valider/Rejeter les offres & Haute \\
    \hline
    Admin & Tableau de bord analytique & Moyenne \\
    \hline
    Finance & Gérer les transactions et virements & Haute \\
    \hline
  \end{tabularx}
\end{table}

\subsection{Besoins non fonctionnels}

Les besoins non fonctionnels définissent les contraintes qualitatives que doit respecter
le système indépendamment de ses fonctionnalités métier.

\begin{table}[H]
  \centering
  \caption{Besoins non fonctionnels du système AfriMarket}
  \label{tab:besoins_non_fonctionnels}
  \begin{tabularx}{\textwidth}{|l|X|X|}
    \hline
    \textbf{Catégorie} & \textbf{Exigence} & \textbf{Indicateur de mesure} \\
    \hline
    Performance & Temps de réponse de l'API inférieur à 500 ms & Mesure avec Postman / JMeter \\
    \hline
    Disponibilité & Uptime du backend supérieur à 99\% & Monitoring uptime \\
    \hline
    Sécurité & Authentification par \gls{jwt} & Rotation du token toutes les 24h \\
    \hline
    Sécurité & Chiffrement des mots de passe (BCrypt) & Audit du code SecurityConfig \\
    \hline
    Scalabilité & Architecture stateless permettant la mise à l'échelle horizontale & Containerisation Docker \\
    \hline
    Utilisabilité & Interface utilisable sans formation préalable & Tests utilisateurs \\
    \hline
    Maintenabilité & Code couvert à 60\% par des tests unitaires & Rapport JaCoCo \\
    \hline
    Compatibilité & Support Android 8+ et iOS 14+ & Tests sur émulateurs \\
    \hline
    Hors-ligne & Affichage des offres récentes sans connexion & Mode dégradé avec cache local \\
    \hline
  \end{tabularx}
\end{table}

% ─── 3.2 Conception ────────────────────────────────────────
\section{Conception}

\subsection{Modèle conceptuel des données}

Le modèle de données d'AfriMarket s'articule autour de neuf entités principales. L'entité
\texttt{User} centralise les informations de tous les acteurs (consommateurs, producteurs,
gestionnaires de relais, administrateurs, agents financiers) grâce à un champ
\texttt{role} de type énuméré. L'entité \texttt{Offer} modélise une offre groupée avec
ses seuils, son statut et ses photos associées. L'entité \texttt{Order} représente la
participation d'un consommateur à une offre, avec capture du prix au moment de l'achat.
Les entités \texttt{Transaction} et \texttt{Wallet} (virtuel) gèrent les flux financiers
internes. L'entité \texttt{RelayPoint} regroupe les points relais géoréférencés associés
à des \texttt{Zone}. Enfin, les entités \texttt{Rating} et \texttt{Dispute} couvrent les
mécanismes de réputation et de résolution des litiges.

% Diagramme entité-relation simplifié en TikZ
\begin{figure}[H]
  \centering
  \begin{tikzpicture}[
      entity/.style={rectangle, draw=NavyBlue, thick, fill=NavyBlue!10,
                     rounded corners=2pt, minimum width=2.4cm, minimum height=0.8cm,
                     font=\small\bfseries},
      rel/.style={diamond, draw=BrickRed, fill=BrickRed!10, thick,
                  aspect=2, font=\scriptsize},
      line/.style={draw, thick, -{Latex}}
  ]
    % Entités
    \node[entity] (user)    at (0,0)      {User};
    \node[entity] (offer)   at (4,0)      {Offer};
    \node[entity] (order)   at (4,-3)     {Order};
    \node[entity] (tx)      at (8,-3)     {Transaction};
    \node[entity] (relay)   at (8,0)      {RelayPoint};
    \node[entity] (zone)    at (8,2.5)    {Zone};
    \node[entity] (photo)   at (4,2.5)    {OfferPhoto};
    \node[entity] (dispute) at (0,-3)     {Dispute};
    \node[entity] (rating)  at (0,-5)     {Rating};

    % Relations
    \draw[line] (user)    -- node[above,font=\scriptsize]{1} (offer)
                           node[midway,above,font=\scriptsize]{produit\ *};
    \draw[line] (user)    -- node[left,font=\scriptsize]{1} (order)
                           node[midway,left,font=\scriptsize]{passe\ *} (order);
    \draw[line] (offer)   -- node[right,font=\scriptsize]{1} (order)
                           node[midway,right,font=\scriptsize]{*};
    \draw[line] (order)   -- node[above,font=\scriptsize]{0..1} (tx)
                           node[midway,above,font=\scriptsize]{*};
    \draw[line] (relay)   -- (zone);
    \draw[line] (offer)   -- (photo);
    \draw[line] (order)   -- (relay);
    \draw[line] (order)   -- (dispute);
    \draw[line] (user)    -- (rating);
  \end{tikzpicture}
  \caption{Diagramme entité-association simplifié d'AfriMarket}
  \label{fig:erd}
\end{figure}

\subsection{Diagramme des cas d'utilisation}

Le diagramme de cas d'utilisation ci-dessous représente les interactions entre les cinq
acteurs principaux et les grandes fonctions du système.

\begin{figure}[H]
  \centering
  \begin{tikzpicture}[
      actor/.style={font=\small, text centered, text width=1.5cm},
      usecase/.style={ellipse, draw=NavyBlue, fill=NavyBlue!8, thick,
                      font=\scriptsize, text width=2.2cm, text centered,
                      minimum height=0.9cm},
      boundary/.style={rectangle, draw=gray, dashed, thick, rounded corners,
                       inner sep=10pt}
  ]
    % Acteurs (gauche)
    \node[actor] (consommateur) at (-5, 0) {
      \tikz\draw[fill=black] (0,0.4) circle (0.15cm);
      \tikz\draw[thick] (0,0) -- (0,-0.7) (-0.3,-0.3) -- (0.3,-0.3) (0,-0.7)--(-0.2,-1.1) (0,-0.7)--(0.2,-1.1);
      \\Consommateur};
    \node[actor] (producteur)   at (-5,-4) {
      \tikz\draw[fill=black] (0,0.4) circle (0.15cm);\\Producteur};
    \node[actor] (admin)        at (-5,-8) {
      \tikz\draw[fill=black] (0,0.4) circle (0.15cm);\\Admin};
    \node[actor] (finance)      at (7, -2) {
      \tikz\draw[fill=black] (0,0.4) circle (0.15cm);\\Finance};
    \node[actor] (relay)        at (7, -6) {
      \tikz\draw[fill=black] (0,0.4) circle (0.15cm);\\Relais};

    % Cas d'utilisation
    \node[usecase] (browsing)   at (0, 1)   {Parcourir les offres};
    \node[usecase] (join)       at (0,-1)   {Rejoindre une offre};
    \node[usecase] (pay)        at (0,-3)   {Payer (MoMo)};
    \node[usecase] (confirm)    at (0,-5)   {Confirmer réception};
    \node[usecase] (create)     at (3, 1)   {Créer offre};
    \node[usecase] (wallet)     at (3,-1)   {Gérer wallet};
    \node[usecase] (validate)   at (3,-5)   {Valider offres};
    \node[usecase] (dashboard)  at (3,-7)   {Tableau de bord};
    \node[usecase] (txmgmt)     at (5.5,-2) {Gérer transactions};

    % Liaisons consommateur
    \draw[thin] (consommateur) -- (browsing);
    \draw[thin] (consommateur) -- (join);
    \draw[thin] (consommateur) -- (pay);
    \draw[thin] (consommateur) -- (confirm);
    % Liaisons producteur
    \draw[thin] (producteur)   -- (create);
    \draw[thin] (producteur)   -- (wallet);
    % Liaisons admin
    \draw[thin] (admin)        -- (validate);
    \draw[thin] (admin)        -- (dashboard);
    % Liaisons finance
    \draw[thin] (finance)      -- (txmgmt);
    \draw[thin] (finance)      -- (wallet);
    % Liaisons relais
    \draw[thin] (relay)        -- (confirm);
  \end{tikzpicture}
  \caption{Diagramme des cas d'utilisation d'AfriMarket}
  \label{fig:usecase}
\end{figure}

\subsection{Cycle de vie d'une offre et d'une commande}

Le cycle de vie d'une offre dans AfriMarket suit une séquence d'états précisément
définis. Une offre débute à l'état \textit{DRAFT} lors de sa création par le producteur,
puis passe à \textit{PENDING} après sa soumission à validation. L'administrateur la valide
(\textit{ACTIVE}) ou la rejette (\textit{REJECTED}). Une offre active collecte des
commandes jusqu'à atteindre son seuil minimal (\textit{THRESHOLD\_REACHED}), après quoi
la livraison est confirmée (\textit{DELIVERING}), puis l'offre est clôturée
(\textit{COMPLETED}) ou annulée (\textit{CANCELLED}) si le seuil n'est pas atteint avant
expiration.

\begin{figure}[H]
  \centering
  \begin{tikzpicture}[
      state/.style={rectangle, draw=NavyBlue, fill=NavyBlue!10, rounded corners=4pt,
                    thick, minimum width=2.8cm, minimum height=0.7cm, font=\small},
      arrow/.style={-{Latex}, thick, NavyBlue},
      label/.style={font=\scriptsize, fill=white, inner sep=1pt}
  ]
    \node[state] (draft)     at (0,0)      {DRAFT};
    \node[state] (pending)   at (3.5,0)    {PENDING};
    \node[state] (active)    at (7,0)      {ACTIVE};
    \node[state] (threshold) at (7,-2)     {THRESHOLD\_REACHED};
    \node[state] (delivering)at (7,-4)     {DELIVERING};
    \node[state] (completed) at (3.5,-4)   {COMPLETED};
    \node[state] (cancelled) at (0,-4)     {CANCELLED};
    \node[state,draw=BrickRed,fill=BrickRed!10] (rejected) at (3.5,-2) {REJECTED};

    \draw[arrow] (draft)     -- node[label,above]{submit} (pending);
    \draw[arrow] (pending)   -- node[label,above]{approve} (active);
    \draw[arrow] (pending)   -- node[label,right]{reject} (rejected);
    \draw[arrow] (active)    -- node[label,right]{seuil atteint} (threshold);
    \draw[arrow] (active)    -- node[label,below,pos=0.3]{expire} (cancelled);
    \draw[arrow] (threshold) -- node[label,right]{livraison} (delivering);
    \draw[arrow] (delivering)-- node[label,above]{complet} (completed);
    \draw[arrow] (active)    to[bend left=20] node[label,above]{annuler} (cancelled);
  \end{tikzpicture}
  \caption{Cycle de vie d'une offre dans AfriMarket}
  \label{fig:offer_lifecycle}
\end{figure}

% ─── 3.3 Architecture ──────────────────────────────────────
\section{Architecture du système}

L'architecture d'AfriMarket suit un schéma en couches (layered architecture) côté backend
et un pattern \gls{mvvm} côté applications mobiles.

\subsection{Architecture globale}

Le backend Spring Boot est organisé en quatre couches distinctes : la couche
\textit{Controller} expose les endpoints \gls{rest} et prend en charge la validation des
entrées ; la couche \textit{Service} contient la logique métier ; la couche
\textit{Repository} abstrait l'accès aux données via Spring Data \gls{jpa} ; enfin, la
couche \textit{Model} définit les entités persistées dans la base de données SQLite via
l'\gls{orm} Hibernate.

\begin{figure}[H]
  \centering
  \begin{tikzpicture}[
      app/.style={rectangle, draw=Teal, fill=Teal!10, thick,
                  rounded corners, minimum width=2.5cm, minimum height=1cm, font=\small},
      layer/.style={rectangle, draw=NavyBlue, fill=NavyBlue!8, thick,
                    minimum width=8cm, minimum height=0.8cm, font=\small\bfseries},
      arrow/.style={-{Latex}, thick}
  ]
    % Applications mobiles
    \node[app] (client) at (-4.5, 3.5) {AfriMarket\\Client};
    \node[app] (prod)   at (-1.5, 3.5) {AfriMarket\\Producteur};
    \node[app] (admin)  at (1.5,  3.5) {Interface\\Admin (Web)};

    % Backend layers
    \node[layer] (ctrl)  at (0, 2) {Controller Layer (REST API / \gls{jwt})};
    \node[layer] (svc)   at (0, 1) {Service Layer (Logique Métier)};
    \node[layer] (repo)  at (0, 0) {Repository Layer (Spring Data JPA)};
    \node[layer] (model) at (0,-1) {Model Layer (Entités JPA)};

    % Base de données
    \node[rectangle, draw=ForestGreen, fill=ForestGreen!10, thick,
          rounded corners, minimum width=3cm, minimum height=0.8cm,
          font=\small] (db) at (0,-2.2) {Base SQLite / afrimarket.db};

    % Flèches
    \draw[arrow] (client) -- (ctrl);
    \draw[arrow] (prod)   -- (ctrl);
    \draw[arrow] (admin)  -- (ctrl);
    \draw[arrow] (ctrl)   -- (svc);
    \draw[arrow] (svc)    -- (repo);
    \draw[arrow] (repo)   -- (model);
    \draw[arrow] (model)  -- (db);
  \end{tikzpicture}
  \caption{Architecture en couches d'AfriMarket}
  \label{fig:architecture}
\end{figure}

\subsection{Organisation des applications mobiles}

Les deux applications mobiles (AfriMarket Client et AfriMarket Producteur) suivent la
même structure dictée par Expo Router. Les routes sont organisées par fichiers dans le
dossier \texttt{app/}, avec des groupes de routes pour l'authentification (\texttt{(auth)})
et les onglets principaux (\texttt{(tabs)}). La couche \texttt{services/} encapsule les
appels \gls{api} vers le backend. La couche \texttt{store/} gère l'état global avec
Zustand, bibliothèque de gestion d'état légère et sans boilerplate. La couche
\texttt{components/} organise les composants réutilisables, eux-mêmes subdivisés par
domaine fonctionnel.

% ─── 3.4 Bilan ─────────────────────────────────────────────
\section{Bilan méthodologique}

La phase d'analyse et de conception a permis d'établir des fondations solides pour le
développement. L'identification claire des cinq rôles utilisateurs et de leurs besoins
spécifiques a guidé la définition des entités du domaine et l'organisation des routes
\gls{api}. Le choix d'une architecture en couches standard, combiné à une séparation nette
entre les applications mobiles et le backend via une \gls{api} \gls{rest}, garantit une
bonne maintenabilité et ouvre la voie à des évolutions futures, comme l'ajout d'un service
de notifications push ou l'intégration d'une \gls{api} de cartographie pour la
géolocalisation des points relais. La décision de reporter la migration vers une base de
données de production (PostgreSQL) à une phase ultérieure, en utilisant SQLite pour le
développement, a permis d'accélérer les cycles d'itération.

% ============================================================
%  CHAPITRE 4 – IMPLÉMENTATION ET RÉSULTATS
% ============================================================
\chapter{Implémentation et résultats}
\markboth{Chapitre 4 -- Implémentation et résultats}{}

% ─── 4.1 Choix technologiques ──────────────────────────────
\section{Choix technologiques}

\subsection{Frontend mobile : React Native avec Expo}

React Native a été retenu pour le développement des applications mobiles AfriMarket
Client et AfriMarket Producteur. Ce framework JavaScript, développé par Meta, permet de
produire des applications natives pour iOS et Android à partir d'une base de code unique,
réduisant ainsi considérablement les délais de développement. Expo, surcouche qui simplifie
la gestion du cycle de vie React Native, apporte le routeur par fichiers (Expo Router),
la gestion simplifiée des assets et des outils de débogage efficaces.
TypeScript a été adopté comme langage de programmation pour bénéficier du typage statique,
réduisant les erreurs à l'exécution et améliorant l'expérience de développement avec
l'autocomplétion \gls{ide} \cite{cherny2019}.

\subsection{Gestion d'état : Zustand}

Zustand a été préféré à Redux pour la gestion d'état global dans les applications mobiles.
Sa philosophie minimaliste (absence de reducers, d'actions et de dispatchers) réduit le
code boilerplate tout en offrant des performances équivalentes. Les stores
(\texttt{auth.store.ts}, \texttt{offer.store.ts}, \texttt{wallet.store.ts}) encapsulent
chacun l'état et les actions d'un domaine fonctionnel.

\subsection{Backend : Spring Boot 4 avec Java 17}

Spring Boot 4, basé sur Jakarta EE 10, constitue le socle du backend. Spring Security
gère l'authentification et les autorisations : un filtre \gls{jwt} personnalisé
(\texttt{JwtAuthFilter}) intercepte chaque requête, valide le token et injecte
l'utilisateur dans le contexte de sécurité. Spring Data \gls{jpa} avec Hibernate assure
la persistance objet-relationnelle. Lombok élimine le code répétitif
(getters/setters/constructeurs). La base de données SQLite, via le driver Xerial, est
utilisée pour le développement ; son fichier \texttt{afrimarket.db} constitue un artefact
deployable autonome.

\subsection{Administration web : Spring MVC avec Thymeleaf}

L'interface d'administration est une application web classique rendue côté serveur via le
moteur de templates Thymeleaf, intégrée au même serveur Spring Boot. Cette approche
simplifie le déploiement (un seul artefact) et exploite Spring Security pour sécuriser
l'accès à l'espace admin par une authentification formulaire distincte de l'authentification
\gls{jwt} utilisée par les applications mobiles.

\begin{table}[H]
  \centering
  \caption{Récapitulatif des technologies utilisées}
  \label{tab:technologies}
  \begin{tabular}{|l|l|l|}
    \hline
    \textbf{Composante} & \textbf{Technologie} & \textbf{Version} \\
    \hline
    Application consommateur & React Native / Expo & SDK 53 \\
    \hline
    Application producteur  & React Native / Expo & SDK 53 \\
    \hline
    Langage frontend & TypeScript & 5.x \\
    \hline
    Gestion d'état & Zustand & 4.x \\
    \hline
    Backend \gls{api} & Spring Boot & 4.0.6 \\
    \hline
    Langage backend & Java & 17 \\
    \hline
    Sécurité & Spring Security + \gls{jwt} & 6.x \\
    \hline
    Persistance & Spring Data JPA / Hibernate & 6.x \\
    \hline
    Base de données & SQLite & 3.x \\
    \hline
    Templates admin & Thymeleaf & 3.x \\
    \hline
    Build backend & Maven & 3.9 \\
    \hline
    Gestion dépendances frontend & npm & 10.x \\
    \hline
  \end{tabular}
\end{table}

% ─── 4.2 Implémentation ────────────────────────────────────
\section{Présentation de l'application AfriMarket}

\subsection{Application consommateur}

L'application consommateur (AfriMarket Client) propose cinq écrans principaux accessibles
via une barre de navigation inférieure. L'écran \textit{Découvrir} affiche les offres
actives filtrables par catégorie et par zone, avec un champ de recherche textuelle. Chaque
carte d'offre présente le titre, la catégorie, le prix unitaire, le producteur et une
barre de progression indiquant le pourcentage du seuil minimal atteint. En appuyant sur
une offre, l'utilisateur accède à l'écran de détail qui récapitule toutes les
informations, liste les points relais disponibles dans la zone et permet de spécifier
la quantité souhaitée avant de rejoindre l'offre. L'écran \textit{Commandes} liste
l'historique des participations avec leur statut. L'écran \textit{Explorer} permet de
parcourir les zones géographiques. Un écran de profil gère les informations personnelles.

\begin{figure}[H]
  \centering
  \begin{tikzpicture}
    % Maquette simplifiée de l'écran Découvrir (fil d'offres)
    \draw[thick, rounded corners=6pt] (0,0) rectangle (5,9);
    % Barre de statut
    \fill[NavyBlue!80] (0,8.5) rectangle (5,9);
    \node[white, font=\scriptsize] at (2.5,8.75) {AfriMarket — Découvrir};
    % Barre de recherche
    \draw[gray!50, rounded corners=3pt, fill=gray!10] (0.2,7.9) rectangle (4.8,8.4);
    \node[gray, font=\scriptsize] at (2.5,8.15) {Rechercher une offre\ldots};
    % Filtres catégorie
    \foreach \x/\lbl in {0.4/Légumes, 1.5/Fruits, 2.6/Tubercules, 3.9/Céréales}{
      \draw[NavyBlue!60, rounded corners=3pt, fill=NavyBlue!15] (\x,7.4) rectangle (\x+0.9,7.75);
      \node[NavyBlue, font=\tiny] at (\x+0.45,7.575) {\lbl};
    }
    % Carte offre 1
    \draw[gray!40, rounded corners=4pt, fill=white, drop shadow] (0.2,5.5) rectangle (4.8,7.2);
    \fill[NavyBlue!20, rounded corners=4pt] (0.2,6.4) rectangle (1.4,7.2);
    \node[font=\tiny\bfseries] at (2.8,7.0) {Tomates fraîches — Yaoundé 4};
    \node[font=\tiny, gray] at (2.8,6.75) {par Jean Nkomo · \(\star\) 4.5};
    \draw[gray!40] (0.5,6.5) rectangle (4.5,6.6);
    \fill[ForestGreen!70] (0.5,6.5) rectangle (3.2,6.6);  % progression 75%
    \node[font=\tiny, ForestGreen] at (2.5,6.35) {75\% — 37,5 kg sur 50 kg min.};
    \node[font=\tiny\bfseries, BrickRed] at (4.0,6.1) {500 FCFA/kg};
    % Carte offre 2
    \draw[gray!40, rounded corners=4pt, fill=white] (0.2,3.7) rectangle (4.8,5.4);
    \fill[NavyBlue!20, rounded corners=4pt] (0.2,4.6) rectangle (1.4,5.4);
    \node[font=\tiny\bfseries] at (2.8,5.2) {Plantains mûrs — Yaoundé 2};
    \node[font=\tiny, gray] at (2.8,4.95) {par Marie Ateba · \(\star\) 4.2};
    \draw[gray!40] (0.5,4.7) rectangle (4.5,4.8);
    \fill[ForestGreen!70] (0.5,4.7) rectangle (2.0,4.8);  % progression 40%
    \node[font=\tiny, ForestGreen] at (2.5,4.55) {40\% — 20 kg sur 50 kg min.};
    \node[font=\tiny\bfseries, BrickRed] at (4.0,4.3) {300 FCFA/kg};
    % Barre navigation
    \fill[gray!15] (0,0) rectangle (5,0.8);
    \foreach \x/\lbl in {0.5/Découvrir, 1.7/Commandes, 2.9/Explorer, 4.1/Profil}{
      \node[font=\tiny, NavyBlue] at (\x+0.3,0.4) {\lbl};
    }
  \end{tikzpicture}
  \caption{Maquette de l'écran « Découvrir » de l'application consommateur}
  \label{fig:screen_discover}
\end{figure}

\subsection{Application producteur}

L'application producteur (AfriMarket Producteur) cible les agriculteurs et leur offre
un tableau de bord synthétique affichant les revenus du mois, le nombre d'offres actives
et le solde du wallet disponible. L'écran de création d'offre guide l'utilisateur à travers
un formulaire multi-étapes : informations générales (titre, description, catégorie), tarification
et quantités (prix unitaire, quantité disponible, seuil minimal, quantité minimale par
acheteur), photos de la production, et dates (disponibilité et expiration). Une fois
soumise, l'offre entre dans la file de validation administrative. L'écran de liste des
offres affiche le statut de chaque offre et, pour les offres actives, la liste des
commandes reçues avec les coordonnées des acheteurs. L'écran Wallet permet au producteur
de visualiser son solde disponible, le montant en attente (commandes confirmées non encore
virées) et l'historique des transactions, et de déclencher une demande de retrait vers
son compte Mobile Money (minimum 500 \gls{fcfa}).

\subsection{Interface d'administration}

L'interface d'administration, accessible via navigateur, est protégée par une
authentification par formulaire. Le tableau de bord principal affiche des indicateurs
clés : chiffre d'affaires total, nombre de commandes actives, utilisateurs enregistrés et
offres en attente de validation. Des vues dédiées permettent de gérer les utilisateurs
(liste, détail, activation/désactivation), les offres (validation, rejet avec motif,
consultation des photos), les commandes (suivi du cycle de vie), les transactions
financières, les points relais (création, modification, association à une zone) et les
paramètres système. L'ensemble est rendu via des templates Thymeleaf avec une mise en
forme CSS administrative sobre et fonctionnelle.

% ─── 4.3 Résultats ─────────────────────────────────────────
\section{Résultats}

\subsection{Fonctionnalités livrées}

L'ensemble des fonctionnalités prioritaires identifiées lors de la phase d'analyse a été
implémenté. Le tableau suivant synthétise l'état de livraison :

\begin{table}[H]
  \centering
  \caption{État de livraison des fonctionnalités}
  \label{tab:livraison}
  \begin{tabularx}{\textwidth}{|X|l|l|}
    \hline
    \textbf{Fonctionnalité} & \textbf{Priorité} & \textbf{Statut} \\
    \hline
    Inscription / Connexion (JWT) & Haute & \textcolor{ForestGreen}{\textbf{Livré}} \\
    \hline
    Parcourir et filtrer les offres & Haute & \textcolor{ForestGreen}{\textbf{Livré}} \\
    \hline
    Rejoindre une offre groupée & Haute & \textcolor{ForestGreen}{\textbf{Livré}} \\
    \hline
    Progression en temps réel & Haute & \textcolor{ForestGreen}{\textbf{Livré}} \\
    \hline
    Paiement Mobile Money & Haute & \textcolor{orange}{\textbf{Simulation}} \\
    \hline
    Créer et soumettre une offre & Haute & \textcolor{ForestGreen}{\textbf{Livré}} \\
    \hline
    Téléversement de photos & Moyenne & \textcolor{ForestGreen}{\textbf{Livré}} \\
    \hline
    Validation admin des offres & Haute & \textcolor{ForestGreen}{\textbf{Livré}} \\
    \hline
    Wallet virtuel et retrait & Haute & \textcolor{ForestGreen}{\textbf{Livré}} \\
    \hline
    Tableau de bord admin & Moyenne & \textcolor{ForestGreen}{\textbf{Livré}} \\
    \hline
    Gestion des points relais & Moyenne & \textcolor{ForestGreen}{\textbf{Livré}} \\
    \hline
    Notifications push & Faible & \textcolor{gray}{\textbf{Non livré}} \\
    \hline
    Géolocalisation des relais & Faible & \textcolor{gray}{\textbf{Non livré}} \\
    \hline
  \end{tabularx}
\end{table}

\subsection{Structure de l'API REST}

L'\gls{api} \gls{rest} expose des ressources organisées selon les conventions du standard.
Les endpoints sont regroupés sous le préfixe \texttt{/api/v1/} et protégés par
l'authentification \gls{jwt}, à l'exception des routes d'authentification
(\texttt{/api/v1/auth/login}, \texttt{/api/v1/auth/register}).

\begin{table}[H]
  \centering
  \caption{Principaux endpoints de l'API AfriMarket}
  \label{tab:api_endpoints}
  \begin{tabular}{|l|l|l|}
    \hline
    \textbf{Méthode} & \textbf{Endpoint} & \textbf{Description} \\
    \hline
    POST & /api/v1/auth/login & Authentification, retourne un JWT \\
    \hline
    POST & /api/v1/auth/register & Inscription d'un nouvel utilisateur \\
    \hline
    GET  & /api/v1/offers & Liste des offres actives (filtrables) \\
    \hline
    GET  & /api/v1/offers/\{id\} & Détail d'une offre \\
    \hline
    POST & /api/v1/offers/\{id\}/join & Rejoindre une offre groupée \\
    \hline
    GET  & /api/v1/offers/mine & Mes offres (producteur) \\
    \hline
    POST & /api/v1/offers & Créer une offre (producteur) \\
    \hline
    POST & /api/v1/offers/\{id\}/submit & Soumettre une offre à validation \\
    \hline
    GET  & /api/v1/orders/mine & Mes commandes (consommateur) \\
    \hline
    GET  & /api/v1/wallet & Solde et infos du wallet \\
    \hline
    POST & /api/v1/wallet/withdraw & Demander un retrait MoMo \\
    \hline
    GET  & /api/v1/zones & Liste des zones géographiques \\
    \hline
    GET  & /api/v1/relays & Liste des points relais \\
    \hline
  \end{tabular}
\end{table}

\subsection{Bilan et perspectives}

Le projet AfriMarket a atteint ses objectifs principaux en livrant une solution mobile
complète et fonctionnelle répondant aux besoins identifiés des producteurs et consommateurs
camerounais. Les trois interfaces (consommateur, producteur, admin) s'articulent
correctement autour du backend unifié, et le mécanisme d'offres groupées constitue la
valeur ajoutée principale du système.

Plusieurs améliorations sont envisagées pour les versions futures. L'intégration réelle
des \gls{api} Mobile Money (MTN et Orange) remplacerait la simulation actuelle et
permettrait un déploiement en production. L'ajout des notifications push (via Firebase
Cloud Messaging) améliorerait l'engagement des utilisateurs. La migration de SQLite vers
PostgreSQL garantirait la fiabilité et les performances en environnement de production.
Enfin, l'implémentation d'un système de recommandation basé sur les préférences d'achat
et la géolocalisation enrichirait l'expérience de découverte des offres.

% ============================================================
%  CONCLUSION GÉNÉRALE
% ============================================================
\chapter*{Conclusion générale}
\addcontentsline{toc}{chapter}{Conclusion générale}

Ce rapport a présenté le projet AfriMarket, une application mobile marketplace agricole
pair-à-pair développée pour le contexte camerounais. Partant d'un problème de terrain bien
identifié — la difficulté pour les agriculteurs péri-urbains d'accéder aux marchés urbains
à prix équitable — nous avons conduit un processus complet allant de l'analyse des besoins
à la livraison d'une solution fonctionnelle.

L'état de l'art a confirmé l'absence de solution adaptée au marché camerounais combinant
offres groupées, réseau de points relais, wallet virtuel et prise en charge native du
Mobile Money. Ce vide justifiait le développement d'AfriMarket comme solution sur mesure
plutôt que l'adaptation d'une plateforme existante.

La phase de conception a permis d'établir un modèle de données robuste autour de neuf
entités interconnectées, un découpage clair des rôles utilisateurs et une architecture en
couches garantissant la maintenabilité. Le choix de React Native avec Expo pour le
frontend et de Spring Boot 4 pour le backend a été motivé par la productivité, la
maturité des écosystèmes et l'adéquation au contexte mobile africain.

Les résultats obtenus attestent de la faisabilité technique de la solution : les
fonctionnalités prioritaires sont livrées et les architectures retenues tiennent leurs
promesses en termes de séparation des préoccupations et d'évolutivité. L'intégration
réelle des paiements Mobile Money et le déploiement sur les stores iOS et Android
constituent les prochaines étapes naturelles pour transformer ce prototype académique en
produit commercialisable.

Au-delà de l'aspect technique, ce projet a été une expérience enrichissante de gestion de
la complexité d'un système multi-acteurs, de prise de décision architecturale sous
contrainte de temps, et d'adaptation des meilleures pratiques du développement logiciel
aux réalités du terrain africain.

% ============================================================
%  ANNEXES
% ============================================================
\appendix
\chapter{Extraits de code significatifs}
\markboth{Annexe A -- Extraits de code}{}

\section{Filtre d'authentification JWT (backend)}

\begin{lstlisting}[language=Java, caption={JwtAuthFilter.java -- Filtre Spring Security},
  label=lst:jwtfilter]
// Filtre Spring Security : intercepte chaque requete HTTP,
// extrait le token JWT de l'en-tete Authorization,
// le valide et injecte l'utilisateur dans le contexte de securite.
@Component
@RequiredArgsConstructor
public class JwtAuthFilter extends OncePerRequestFilter {

    private final JwtService jwtService;
    private final UserDetailsService userDetailsService;

    @Override
    protected void doFilterInternal(HttpServletRequest request,
            HttpServletResponse response, FilterChain filterChain)
            throws ServletException, IOException {

        final String authHeader = request.getHeader("Authorization");

        // Si l'en-tete est absent ou ne commence pas par "Bearer ", on passe
        if (authHeader == null || !authHeader.startsWith("Bearer ")) {
            filterChain.doFilter(request, response);
            return;
        }

        // Extraction du token (on retire le prefixe "Bearer ")
        final String jwt = authHeader.substring(7);
        final String userPhone = jwtService.extractUsername(jwt);

        // Si le token est valide et qu'aucune authentification n'est en cours
        if (userPhone != null &&
            SecurityContextHolder.getContext().getAuthentication() == null) {

            UserDetails userDetails =
                userDetailsService.loadUserByUsername(userPhone);

            if (jwtService.isTokenValid(jwt, userDetails)) {
                // Injection de l'utilisateur dans le contexte de securite
                UsernamePasswordAuthenticationToken authToken =
                    new UsernamePasswordAuthenticationToken(
                        userDetails, null,
                        userDetails.getAuthorities());
                authToken.setDetails(
                    new WebAuthenticationDetailsSource()
                        .buildDetails(request));
                SecurityContextHolder.getContext()
                    .setAuthentication(authToken);
            }
        }
        filterChain.doFilter(request, response);
    }
}
\end{lstlisting}

\section{Store Wallet avec Zustand (frontend producteur)}

\begin{lstlisting}[language=Java, caption={wallet.store.ts -- Gestion d'état avec Zustand},
  label=lst:walletstore]
// Store Zustand pour le wallet du producteur.
// Encapsule l'etat (wallet, transactions, loading) et les actions.
interface WalletState {
  wallet: WalletDto | null;
  transactions: TransactionDto[];
  loading: boolean;
  fetchWallet: () => Promise<void>;
  fetchTransactions: () => Promise<void>;
}

export const useWalletStore = create<WalletState>((set) => ({
  wallet: null,
  transactions: [],
  loading: false,

  // Charge le solde du wallet depuis l'API backend
  fetchWallet: async () => {
    set({ loading: true });
    try {
      const data = await apiGet<WalletDto>(API.WALLET);
      set({ wallet: data });
    } finally {
      set({ loading: false });
    }
  },

  // Charge l'historique des transactions
  fetchTransactions: async () => {
    const data = await apiGet<TransactionDto[]>(API.TRANSACTIONS);
    set({ transactions: data });
  },
}));
\end{lstlisting}

\chapter{Modèle de données complet}
\markboth{Annexe B -- Modèle de données}{}

\begin{table}[H]
  \centering
  \caption{Entités de la base de données AfriMarket}
  \label{tab:entities}
  \begin{tabularx}{\textwidth}{|l|X|l|}
    \hline
    \textbf{Entité} & \textbf{Champs principaux} & \textbf{Relations} \\
    \hline
    User & id, phone, fullName, role, status, momoNumber, momoProvider, zone, ratingAvg & ManyToOne Zone \\
    \hline
    Offer & id, producer, title, category, unit, pricePerUnit, availableQty, minThreshold, minQtyPerBuyer, currentQtyOrdered, status, expiresAt & ManyToOne User, OneToMany OfferPhoto \\
    \hline
    Order & id, offer, buyer, relay, qtyOrdered, unitPriceSnapshot, totalAmount, status, paymentProvider, paidAt & ManyToOne Offer, User, RelayPoint \\
    \hline
    Transaction & id, order, type, amount, commission, netAmount, momoRef, status & ManyToOne Order \\
    \hline
    RelayPoint & id, name, address, zone, latitude, longitude & ManyToOne Zone \\
    \hline
    Zone & id, name, region & -- \\
    \hline
    OfferPhoto & id, offer, url, isPrimary & ManyToOne Offer \\
    \hline
    Dispute & id, order, reason, status, resolution & ManyToOne Order \\
    \hline
    Rating & id, rater, rated, offer, score, comment & ManyToOne User, Offer \\
    \hline
  \end{tabularx}
\end{table}

% ============================================================
%  BIBLIOGRAPHIE (auto via biblatex)
% ============================================================
\printbibliography[heading=bibintoc, title={Références bibliographiques}]

\end{document}
